\documentclass[11pt,a4paper]{article}

\usepackage[margin=1in]{geometry}
\usepackage{lmodern}
\usepackage[T1]{fontenc}
\usepackage[utf8]{inputenc}
\usepackage{microtype}
\usepackage{amsmath,amssymb,amsthm,mathtools}
\usepackage{array}
\usepackage{booktabs}
\usepackage{enumitem}
\usepackage[hidelinks]{hyperref}

\setlist[itemize]{leftmargin=1.4em,itemsep=0.25em,topsep=0.35em}
\setlist[enumerate]{leftmargin=1.6em,itemsep=0.25em,topsep=0.35em}

\theoremstyle{definition}
\newtheorem{definition}{Definition}
\newtheorem{proposition}{Proposition}
\newtheorem{remark}{Remark}

\newcommand{\ket}[1]{\lvert #1\rangle}
\newcommand{\bra}[1]{\langle #1\rvert}
\newcommand{\braket}[2]{\langle #1\vert #2\rangle}
\newcommand{\abs}[1]{\lvert #1\rvert}
\newcommand{\norm}[1]{\lVert #1\rVert}
\newcommand{\Tr}{\operatorname{Tr}}
\newcommand{\dd}{\mathrm{d}}
\newcommand{\ii}{\mathrm{i}}
\newcommand{\eps}{\varepsilon}
\newcommand{\order}{\mathcal{O}}

\title{Incorporating DRAG into Space-Curve Quantum Control\\
and Leakage-Aware BARQ}
\author{Technical design note for the \texttt{qurveros} package}
\date{23 October 2026}

\begin{document}
\maketitle

\begin{abstract}
This note derives how Derivative Removal by Adiabatic Gate (DRAG) can be
combined with space-curve quantum control (SCQC) and the B\'ezier ansatz for
robust quantum control (BARQ), as implemented in \texttt{qurveros}.  The
central result is that a DRAG correction applied after a completed BARQ design
is inexpensive, but it does not automatically preserve the exact geometric
robustness conditions.  It preserves the computational gate only
perturbatively in \(1/(\alpha T)\) and generally introduces mixed errors of
order \((\delta_z T)^2/(\alpha T)^2\).  A leakage-aware BARQ objective can
instead impose four real first-order leakage constraints while recomputing
dephasing and amplitude sensitivities under the corrected qutrit evolution.
The note also records the associated boundary conditions, objective functions,
spectral and smoothness extensions, the relation to a perturbative
\(\mathbb{C}^2\) leakage curve, and the scaling that can reasonably be
expected before a numerical campaign.
\end{abstract}

\tableofcontents

\section{Scope and conventions}

SCQC and BARQ design a single-qubit gate by constructing a space curve whose
curvature and torsion determine the drive.  Those constructions are exact for
a two-level Hamiltonian.  A transmon, or any weakly anharmonic ladder, has at
least one additional level.  DRAG suppresses the resulting leakage by a
near-identity frame transformation whose generator is proportional to the
drive~\cite{Motzoi2009}.  The two ideas therefore meet at the pulse, not
inside the original three-dimensional Frenet frame.

Throughout, \(P=\ket0\bra0+\ket1\bra1\) projects onto the computational
subspace and \(Q=I-P\).  Time \(t\) runs from \(0\) to the gate duration \(T\).
The dimensionless time used by \texttt{qurveros} is
\begin{equation}
  s=\frac{t}{T}\in[0,1].
\end{equation}
A dot denotes \(\dd/\dd t\) and a prime denotes \(\dd/\dd s\).

The minimal physical model is a driven three-level ladder.  In the frame
rotating at the carrier frequency, and after the rotating-wave approximation,
\begin{equation}
\label{eq:h3}
\begin{aligned}
  \frac{H_3(t)}{\hbar}
  &= \delta(t)\ket1\bra1+(2\delta(t)+\alpha)\ket2\bra2 \\
  &\quad+\frac12\Bigl\{\eps(t)\bigl(\ket0\bra1+\lambda\ket1\bra2\bigr)
      +\text{h.c.}\Bigr\}.
\end{aligned}
\end{equation}
The complex envelope is \(\eps(t)=\Omega_x(t)+\ii\Omega_y(t)\), \(\lambda\) is
the ratio of the \(1\leftrightarrow2\) and \(0\leftrightarrow1\) couplings, and
\begin{equation}
  \alpha=\omega_{12}-\omega_{01}
\end{equation}
is the anharmonicity.  For a transmon one typically has \(\lambda=\sqrt2\) and
\(\alpha<0\).  Every quadrature sign below is tied to this convention.  Taking
the opposite complex-envelope convention, or replacing \(\alpha\) by
\(\abs\alpha\), flips signs but not the geometric conditions.

The \texttt{qurveros} control dictionary stores dimensionless fields.  With
\(u(s)\) the complex envelope and \(d(s)\) the detuning,
\begin{equation}
\label{eq:normalized}
  u(s)=T\eps(Ts),\qquad d(s)=T\delta(Ts),\qquad \nu=\alpha T,\qquad
  \eta=\frac1{\abs\nu}.
\end{equation}
In the SCQC mapping, up to the package's sign conventions,
\(\omega(s)=T\kappa(t)\) is the normalized curvature, the phase or detuning is
fixed by the torsion, and \(T\) is the length of the curve.  DRAG is an
asymptotic expansion in \(\eta\).  It is appropriate when the envelope and its
phase vary slowly relative to \(\abs\alpha\), and it becomes progressively less
reliable when \(\abs{\alpha}T\) is only a few radians.

\begin{remark}[What is derived and what is a proxy]
Sections~\ref{sec:leakage}, \ref{sec:drag}, and~\ref{sec:robust} distinguish
three objects: an exact spectral identity for a bare envelope, a first-order
Dyson formula for gate leakage, and the original fixed-axis DRAG
transformation.  A complex-envelope formula for a rapidly varying SCQC phase is
the natural extension of those results, but it remains an ansatz until it is
checked by a qutrit Schrieffer--Wolff expansion or by direct simulation.
\end{remark}

\section{First-order leakage}
\label{sec:leakage}

Let \(U_q(t)\) be the ideal evolution generated by the computational part of
\eqref{eq:h3}, embedded in the qutrit space, and let \(\tilde U_q(s)=U_q(Ts)\).
Treat the \(1\leftrightarrow2\) coupling as a perturbation.  For a computational
input \(\ket j\), the first-order amplitude to occupy \(\ket2\) at time \(T\) is
\begin{equation}
\label{eq:dyson}
  A^{(1)}_{2j}(T)
  =-\ii\frac{\lambda}{2}\int_0^T\dd t\;
    \eps^*(t)\,e^{\ii\alpha t}\,\braket{1}{U_q(t)j}.
\end{equation}
Changing to normalized variables cancels one factor of \(T\):
\begin{equation}
  A^{(1)}_{2j}
  =-\ii\frac{\lambda}{2}\int_0^1\dd s\;
    u^*(s)\,e^{\ii\nu s}\,\braket{1}{\tilde U_q(s)j}.
\end{equation}
Averaging the resulting probability over the two computational inputs gives
\begin{equation}
\label{eq:leak-prob}
  \bar L^{(1)}
  =\frac12\sum_{j=0,1}\abs{A^{(1)}_{2j}}^2
  =\frac{\lambda^2}{8}\sum_{j=0,1}
    \abs{\ell_j(\nu)}^2,
\end{equation}
where the leakage moments are
\begin{equation}
\label{eq:ell}
  \ell_j(\nu)
  \equiv\int_0^1\dd s\;
  u^*(s)\,e^{\ii\nu s}\,\braket{1}{\tilde U_q(s)j}.
\end{equation}

\begin{proposition}[Nominal first-order leakage closure]
\label{prop:leakage}
The first-order leakage amplitude vanishes for every computational input if and
only if
\begin{equation}
\label{eq:leak-conditions}
  \ell_0(\nu)=0,\qquad \ell_1(\nu)=0.
\end{equation}
These are two complex equations, hence four independent real constraints.
\end{proposition}

The cheaper proxy \(\int_0^1 u(s)e^{\ii\nu s}\,\dd s=0\) forces a notch in the
bare envelope spectrum.  It is not equivalent to \eqref{eq:leak-conditions},
because the integrands contain the evolving factor
\(\braket{1}{\tilde U_q(s)j}\).  That factor is available from the ideal qubit
propagator already computed by SCQC, so the stronger condition does not require
a qutrit simulation.

If a physical parameter \(q\) is uncertain, first-order insensitivity of the
leakage cancellation adds
\begin{equation}
  \partial_q\ell_0(\nu)=0,\qquad \partial_q\ell_1(\nu)=0,
\end{equation}
another four real constraints for each such parameter.  The important choices
are anharmonicity, a common amplitude scale, and a static detuning.

\section{DRAG as a post-processor}
\label{sec:drag}

\subsection{Fixed-axis transformation}

Suppose first that the intended envelope \(E_x(t)\) is real and a second
quadrature \(E_y(t)\) is available for correction.  Following
Motzoi~et~al.~\cite{Motzoi2009}, remove the direct leakage coupling by the
near-identity transformation
\begin{equation}
\label{eq:V}
  V(t)=\exp\left\{-\ii\frac{E_x(t)}{2\alpha}
    \bigl(Y_{01}+\lambda Y_{12}\bigr)\right\}.
\end{equation}
The adiabaticity parameter is \(E_x/\alpha\).  To first order in that
parameter, the transformed Hamiltonian contains the inertial term
\begin{equation}
  \frac12\left(E_y+\frac{\dot E_x}{\alpha}\right)(Y_{01}+\lambda Y_{12})
\end{equation}
and an ac-Stark shift of \(\ket1\) proportional to
\begin{equation}
  \delta-\frac{(\lambda^2-4)E_x^2}{4\alpha}.
\end{equation}
Thus the first-order DRAG controls are
\begin{equation}
\label{eq:drag-fixed}
  E_y(t)=-\frac{\dot E_x(t)}{\alpha},
  \qquad
  \delta(t)=\frac{(\lambda^2-4)E_x(t)^2}{4\alpha}.
\end{equation}
For a transmon, \(\lambda^2=2\), so the detuning simplifies to
\(\delta\approx -E_x^2/(2\alpha)\).  Higher-order terms in
Ref.~\cite{Motzoi2009} correct the in-phase envelope, the derivative
quadrature, and the detuning at orders \(E_x^3/\alpha^2\), \(E_x^2\dot
E_x/\alpha^3\), and \(E_x^4/\alpha^3\).  They are useful once \(\eta\) is no
longer small, but they are not required for the first-order theory.

The frame transformation has no net effect on the gate precisely when
\begin{equation}
  V(0)=V(T)=I.
\end{equation}
For \eqref{eq:V}, this is the endpoint condition \(E_x(0)=E_x(T)=0\).

\subsection{Complex envelope}

A general SCQC pulse already uses both quadratures.  The covariant extension of
\eqref{eq:drag-fixed} is
\begin{equation}
\label{eq:drag-complex}
  \eps_D(t)=\eps(t)-\ii\frac{\beta}{\alpha}\dot\eps(t),
\end{equation}
with a dimensionless calibration parameter \(\beta\) equal to one in the
analytic first-order theory.  In normalized variables,
\begin{equation}
\label{eq:drag-u}
  u_D(s)=u(s)-\ii\frac{\beta}{\nu}u'(s).
\end{equation}
Equivalently, in Cartesian components,
\begin{equation}
\label{eq:drag-cartesian}
  \omega_{x,D}=\omega_x+\frac{\beta}{\nu}\omega_y',
  \qquad
  \omega_{y,D}=\omega_y-\frac{\beta}{\nu}\omega_x'.
\end{equation}
If \(u=\omega e^{\ii\phi}\), the same correction is
\begin{equation}
\label{eq:drag-polar}
  u_D(s)=e^{\ii\phi(s)}
  \left[\omega(s)+\frac{\beta}{\nu}\omega(s)\phi'(s)
        -\ii\frac{\beta}{\nu}\omega'(s)\right].
\end{equation}
The term proportional to \(\omega\phi'\) is a radial correction generated by
phase modulation.  It is absent from textbook fixed-axis DRAG and should not
be dropped for an SCQC pulse.  Derivatives should be taken on the Cartesian
quadratures: near a signed-envelope zero, \(\omega\) and \(\phi\) separately
can have spurious singularities even though \(\omega_x\) and \(\omega_y\) are
smooth.

\begin{proposition}[Bare spectral notch]
\label{prop:notch}
Suppose \(\eps(0)=\eps(T)=0\).  Then
\begin{equation}
  \int_0^T\left(\eps(t)-\ii\frac{\dot\eps(t)}{\alpha}\right)e^{\ii\alpha t}\,\dd t
  =0.
\end{equation}
\end{proposition}

\begin{proof}
Integrate the derivative term by parts.  The boundary term is
\((\eps(T)e^{\ii\alpha T}-\eps(0))/\alpha=0\), while
\(-\ii\alpha\int\eps e^{\ii\alpha t}\,\dd t\) cancels the remaining integral of
\(\eps e^{\ii\alpha t}\).
\end{proof}

Proposition~\ref{prop:notch} does not imply
Proposition~\ref{prop:leakage}.  The latter weights \(\eps^*(t)\) by
\(\braket{1}{U_q(t)j}\), and DRAG also changes \(U_q\) at order \(\eta\).
Equation~\eqref{eq:drag-complex} should therefore be read as the appropriate
post-processing ansatz, followed by a small calibration of \(\beta\), an
overall amplitude, and a constant detuning.

\subsection{Boundary conditions and calibration}

The following conditions are listed in the order they become necessary.

\begin{enumerate}
  \item \(u(0)=u(1)=0\) makes the first-order DRAG frame an identity at both
        endpoints.  BARQ already encourages this through its vanishing-envelope
        construction, but a loaded or rescaled pulse should be checked.
  \item \(u'(0)=u'(1)=0\) makes the added derivative quadrature start and end
        continuously.
  \item If corrections through the \(m\)th derivative are used, the derivatives
        of orders \(0,\ldots,m\) should vanish at both endpoints.  Otherwise
        integration by parts leaves boundary kicks in the leakage amplitude.
  \item After constructing \(u_D\), at least the global amplitude and a
        constant phase or detuning offset should be recalibrated.  The
        derivative term does not cancel every Stark and higher-order coherent
        error.
\end{enumerate}

A practical post-processing sequence is:

\begin{enumerate}
  \item Design \(u(s)\) and \(d(s)\) with ordinary BARQ and record the original
        closure and tantrix-area diagnostics.
  \item Apply \eqref{eq:drag-u} and the Stark correction from
        \eqref{eq:drag-fixed}, transformed into normalized units.
  \item On a qutrit model, tune \(\beta\), the amplitude scale, and a constant
        detuning to minimize \(1-F_{\mathrm{avg}}\) and the leakage \(L\)
        defined in Section~\ref{sec:metrics}.
  \item Recompute noise generators with the corrected propagator, rather than
        reusing the original curve tests.
  \item Repeat under variation of \(\alpha\), gain, timing, sampling, and the
        control-line transfer function.
\end{enumerate}

\section{Robustness after DRAG}
\label{sec:robust}

\subsection{Generator conditions}

The original geometric tests are statements about a qubit propagator and a
qubit noise operator.  Both change once the pulse is corrected and embedded in
a qutrit.  Let \(U_D(t)\) be the ideal qutrit evolution generated by the
corrected controls.  For a parameter \(q_a\) with Hamiltonian derivative
\(V_a(t)=\partial H/\partial q_a\), the first Magnus generator is
\begin{equation}
\label{eq:generator}
  G_a=\int_0^T\dd t\; U_D^\dagger(t)\,V_a(t)\,U_D(t).
\end{equation}
Its computational and leakage blocks impose different requirements.

\begin{definition}[First-order qutrit robustness]
The corrected gate is first-order robust against \(q_a\) when
\begin{equation}
\label{eq:robust-blocks}
  PG_aP-\frac12\Tr(PG_aP)\,P=0,
  \qquad QG_aP=0.
\end{equation}
The first equation removes the traceless logical rotation.  The second
prevents the noise from creating leakage to first order.
\end{definition}

For a genuine qubit model one has \(P=I\), and the first equation reduces to
the usual closure condition.  Applying the old closure or tantrix test to the
uncorrected curve is not sufficient, because neither \(U_D\) nor \(V_a\) is the
operator assumed by that test.

\subsection{Perturbation of dephasing closure}

Let \(U_0\) be the computational evolution before the DRAG correction and write
\begin{equation}
  U_D=U_0\exp\bigl[-\ii\eta K_1+\order(\eta^2)\bigr].
\end{equation}
For quasistatic detuning noise, take the interaction-picture operator
\(n_0(t)=U_0^\dagger(t)\sigma_z U_0(t)\) and the associated error curve
\(r_0(T)=\int_0^T n_0(t)\,\dd t\).  Expanding \eqref{eq:generator} gives
\begin{equation}
\label{eq:perturbed-curve}
  r_D(T)=r_0(T)+\ii\eta\int_0^T[K_1(t),n_0(t)]\,\dd t+\order(\eta^2).
\end{equation}
Ordinary SCQC can enforce \(r_0(T)=0\).  It does not enforce the commutator
integral.  Consequently, a pulse that is first-order dephasing-robust before
post-processing generically acquires a residual generator of order
\(\eta\).

For a qubit, a vanishing first-order generator changes the leading coherent
infidelity from quadratic to quartic in the noise strength.  After
unconstrained post-processing the scaling is instead
\begin{equation}
\label{eq:scaling-robust}
\begin{aligned}
  \text{original closure:}\quad
    &1-F=\order\bigl((\delta_zT)^4\bigr),\\
  \text{post-processed DRAG:}\quad
    &1-F=\order\bigl((\delta_zT)^2\eta^2\bigr)+\order\bigl((\delta_zT)^4\bigr).
\end{aligned}
\end{equation}
The mixed term dominates the original quartic protection whenever
\(\abs{\delta_zT}\lesssim\eta\).  This is a scaling statement, not a claim that
the numerical coefficient is large.  Restoring the quartic regime requires the
additional vector condition
\begin{equation}
\label{eq:mixed}
  \int_0^T[K_1(t),n_0(t)]\,\dd t=0,
\end{equation}
or, more directly, the block conditions \eqref{eq:robust-blocks} evaluated on
the final pulse.

\subsection{Amplitude, DRAG gain, and anharmonicity}

If one calibration error \(a\) multiplies both applied quadratures, its
derivative is built from the corrected envelope,
\begin{equation}
  V_a(t)=\frac12\bigl(\Omega_{x,D}(t)X+\Omega_{y,D}(t)Y\bigr),
\end{equation}
and robustness is \(G_a\propto I\) on the computational block.  The original
tantrix-area condition is the special case of this statement before DRAG and
before the qutrit embedding.  If the derivative channel is calibrated
independently, \(\beta\) must be treated as its own parameter and \(G_\beta\)
must vanish separately.  Anharmonicity uncertainty is qualitatively different:
it moves the leakage phase \(e^{\ii\nu s}\) and the coefficient \(1/\nu\) in
\eqref{eq:drag-u} simultaneously.  A first-order robust design therefore
penalizes \(\partial_\alpha M\), where \(M=PU_DP\) is the final computational
block.

Second-order cancellation follows from the next Magnus term.  For two error
channels,
\begin{equation}
  G_{ab}^{(2)}=-\frac{\ii}{2}\int_0^T\dd t_1\int_0^{t_1}\dd t_2\;
  \bigl[V_a^I(t_1),V_b^I(t_2)\bigr],
\end{equation}
where the superscript \(I\) denotes the interaction picture generated by
\(U_D\).  The computational block of \(G_{ab}^{(2)}\), with its trace removed,
is the appropriate second-order condition.  In the original qubit dephasing
model this is the geometric zero-area constraint.

\begin{table}[ht]
\centering
\caption{Noise channels that should be recomputed after DRAG.}
\label{tab:channels}
\begin{tabular}{@{}lll@{}}
\toprule
Uncertainty & Generator & New risk \\
\midrule
Static detuning & \(\partial H/\partial\delta\) & Closure shifted at \(\order(\eta)\) \\
Common I/Q gain & \(\partial H/\partial a\) & Derivative terms scale with the gain \\
DRAG coefficient & \(\partial H/\partial\beta\) & Separate derivative-channel error \\
Anharmonicity & \(\partial H/\partial\alpha\) & Notch frequency and \(1/\alpha\) both move \\
Transfer function & filtered functional derivative & Differentiation amplifies distortion \\
\bottomrule
\end{tabular}
\end{table}

\section{Leakage-aware BARQ}
\label{sec:barq}

Post-processing takes a curve designed for a qubit and changes its pulse.
Leakage-aware BARQ instead varies the free B\'ezier points, and optionally a
few gate-fixing parameters, under a qutrit objective.  Three nested models are
useful.

\subsection{Spectral and perturbative objectives}

The least expensive differentiable objective uses
Proposition~\ref{prop:leakage}:
\begin{equation}
  J_{\mathrm{spec}}=\sum_{j=0,1}\abs{\ell_j(\nu)}^2.
\end{equation}
The bare-envelope proxy
\(J_{\mathrm{bare}}=\abs{\int_0^1 u(s)e^{\ii\nu s}\,\dd s}^2\) is cheaper but
weaker.  The corresponding leakage probability from \eqref{eq:leak-prob} is
\begin{equation}
  J_{\mathrm{leak},1}=\frac{\lambda^2}{8}\bigl(\abs{\ell_0}^2+\abs{\ell_1}^2\bigr).
\end{equation}
Both can be evaluated from the normalized envelope and the ideal qubit
propagator.  Their gradients exist once the cumulative phase and arclength are
computed by a differentiable quadrature rather than by a black-box integrator.

\subsection{Exact qutrit objective}
\label{sec:metrics}

Let \(U\) be the full qutrit propagator and \(M=PUP\) its \(2\times2\)
computational block.  For a pure input drawn uniformly from the qubit Bloch
sphere, the average probability of leaving that subspace is
\begin{equation}
\label{eq:L}
  L=1-\frac12\Tr(M^\dagger M).
\end{equation}
If \(U_g\) is the target qubit unitary, the Haar-averaged state fidelity is
\begin{equation}
\label{eq:F}
  F_{\mathrm{avg}}
  =\frac{\Tr(M^\dagger M)+\abs{\Tr(U_g^\dagger M)}^2}{6}.
\end{equation}
Equation~\eqref{eq:F} follows from
\begin{equation}
  \int\abs{\bra{\psi}A\ket{\psi}}^2\,\dd\psi
  =\frac{\abs{\Tr A}^2+\Tr(A^\dagger A)}{d(d+1)}
\end{equation}
with \(d=2\) and \(A=U_g^\dagger M\).  The leakage term \eqref{eq:L} is already
present in \(F_{\mathrm{avg}}\), but retaining it separately allows leakage to
be weighted more heavily than a coherent logical rotation.  Both quantities
remain meaningful when \(M\) is not unitary.

\subsection{What remains analytic}

BARQ's point-gate-fixing construction fixes the ideal qubit operation and the
endpoint contribution to dephasing closure before numerical optimization.  On
a qutrit, Stark shifts and virtual transitions modify \(M\), so those
constraints no longer guarantee the physical target exactly.  The useful hybrid
is:

\begin{itemize}
  \item retain point-gate fixing, so the zeroth-order qubit gate and curve
        closure remain exact;
  \item constrain enough endpoint control points that \(u\) and \(u'\) vanish
        at \(s=0,1\);
  \item optimize interior points against \(F_{\mathrm{avg}}\), \(L\), and the
        corrected generators \eqref{eq:robust-blocks};
  \item release a global amplitude and a constant detuning during the final
        stage, since qutrit gate fixing is no longer purely analytic;
  \item test four- and five-level truncations for short gates, where
        two-photon and higher-state transitions are outside \eqref{eq:h3}.
\end{itemize}

A concrete staged objective is
\begin{equation}
\label{eq:J}
\begin{aligned}
  J&=w_g(1-F_{\mathrm{avg}})+w_L L
    +w_z\norm{PG_zP-\tfrac12\Tr(PG_zP)P}_F^2
    +w_{zL}\norm{QG_zP}_F^2 \\
   &\quad+w_a\norm{PG_aP-\tfrac12\Tr(PG_aP)P}_F^2
    +w_{aL}\norm{QG_aP}_F^2
    +w_\alpha\norm{\partial_\alpha M}_F^2 \\
   &\quad+w_b\int_0^1\abs{u'(s)}^2\,\dd s
    +w_{bb}\int_0^1\abs{u''(s)}^2\,\dd s
    +w_{\max}\max_s\abs{u(s)} \\
   &\quad+w_{\mathrm{end}}\sum_{k=0,1}\sum_{s_\star\in\{0,1\}}\abs{u^{(k)}(s_\star)}^2.
\end{aligned}
\end{equation}
The weights should be introduced in stages: target and original robustness,
then nominal leakage, then parameter uncertainty, and finally bandwidth and
amplitude.  Turning every term on simultaneously makes the trade-off difficult
to diagnose.

\section{Further extensions}

\subsection{Choosing the gate duration}

For a fixed normalized BARQ waveform, the first-order leakage depends on
duration only through \(\nu=\alpha T\).  A coherent-error search can therefore
minimize \(L_1(\nu)=J_{\mathrm{leak},1}\) over \(T\), subject to the hardware
bound on peak amplitude.  This is stronger than placing a zero of the bare
spectrum at \(\alpha\), because \(\ell_j\) contains \(\tilde U_q\).  Decoherence
selects a finite optimum of a model such as
\begin{equation}
  E(T)\approx L_{\mathrm{leak}}(\alpha T)+\frac{c_1T}{T_1}+\frac{c_\phi T}{T_\phi}
  +E_{\mathrm{control}}(T).
\end{equation}
The coefficients \(c_1\) and \(c_\phi\) depend on how the target gate and noise
channel average relaxation and pure dephasing; they should come from the
open-system simulation rather than from a universal rule.

\subsection{Endpoint smoothness}

Let \(I(\nu)=\int_0^1 g(s)e^{\ii\nu s}\,\dd s\).  If \(g\) and its first
\(m-1\) derivatives vanish at both endpoints, repeated integration by parts
yields
\begin{equation}
  I(\nu)=\left(-\frac1{\ii\nu}\right)^m\int_0^1 g^{(m)}(s)e^{\ii\nu s}\,\dd s,
\end{equation}
and therefore \(\abs I=\order(\abs\nu^{-m})\).  For leakage,
\(g(s)=u^*(s)\braket{1}{\tilde U_q(s)j}\).  Making only \(u\) smooth does not
automatically make every derivative of this product vanish, but it improves the
decay.  The quadratic penalties on \(u'\) and \(u''\) in \eqref{eq:J} suppress
bandwidth even when an endpoint derivative cannot be constrained exactly.

\subsection{Higher-derivative and multi-transition DRAG}

A linear differential filter can place spectral zeros at several unwanted
frequencies \(\Delta_k\):
\begin{equation}
\label{eq:product}
  u_{\mathrm{corr}}(t)
  =\prod_k\left(1-\frac{\ii}{\Delta_k}\frac{\dd}{\dd t}\right)u_0(t).
\end{equation}
Under the Fourier convention \(\int u(t)e^{\ii\omega t}\,\dd t\), each factor
contributes \(1-\omega/\Delta_k\).  Expanding \eqref{eq:product} produces the
higher-derivative DRAG family.  The same construction suppresses several
internal or spectator transitions, provided the corresponding endpoint
derivatives vanish and the coefficients are recalibrated on the full
propagator.  At very short duration, direct and multiphoton transitions beyond
\(\ket2\) can dominate the residual error, so the filter frequencies must
include those transitions or the simulation must retain more levels.

\subsection{Fourier notches and Wah-Wah modulation}

The derivative expansion is not the only way to impose the linear constraints
\(\ell_j(\nu_k)=0\).  A finite spectral ansatz
\(u(s)=\sum_n c_n f_n(s)\) converts each notch into a linear condition on the
coefficients \(c_n\).  Those conditions can be combined with the BARQ gate
area, endpoint, and robustness constraints.  Wah-Wah control is the
crowded-spectrum version of this idea: sideband modulation places a notch near
a neighbour's leakage transition, while a DRAG quadrature suppresses the
addressed qubit's own \(1\leftrightarrow2\)
transition~\cite{Schutjens2013,Vesterinen2014}.  It is relevant when a
spectator transition is closer to the drive than the addressed qubit's
anharmonicity.

\subsection{Relation to multilevel SCQC}

The generalized geometric formalism represents the interaction-picture error
operator by coordinates
\begin{equation}
  q_a(t)=\frac12\Tr\bigl[\Lambda_a U^\dagger(t)QU(t)\bigr],
\end{equation}
where \(\{\Lambda_a\}\) is a generator basis normalized by
\(\Tr(\Lambda_a\Lambda_b)=2\delta_{ab}\)~\cite{Buterakos2021,BarnesReview}.
Leading-order cancellation means that every coordinate curve closes,
\(\int_0^T q_a(t)\,\dd t=0\).  A generic qutrit error gives an eight-dimensional
curve.  Generalized curvatures can be connected recursively to the controls,
but a transmon does not supply enough independent controls to realize an
arbitrary curve: two quadratures and a detuning act with linked
\(0\leftrightarrow1\) and \(1\leftrightarrow2\) matrix elements.

The perturbative intermediate theory avoids that inverse problem.  Leakage
lives in the four-real-dimensional space \(Q\mathcal H P\cong\mathbb{C}^2\),
and the two components of its integrated coordinate are exactly
\((\ell_0,\ell_1)\).  Closure of this leakage curve is
Proposition~\ref{prop:leakage}.  Thus leakage-aware BARQ adds a second,
four-dimensional closure condition without abandoning the existing
\(\mathrm{SU}(2)\) curve used to design the computational gate.

\section{Expected improvement and required evidence}

For a fixed rotation angle, the drive amplitude scales as \(1/T\).  The
dimensionless expansion parameter is therefore \(\eta\sim 1/(\abs\alpha T)\),
up to a pulse-shape coefficient.  A generic off-resonant leakage amplitude
begins at \(\order(\eta)\), so its probability begins at \(\order(\eta^2)\).
First-order DRAG is constructed to cancel that amplitude.  The residual
coherent leakage is then commonly \(\order(\eta^2)\) in amplitude and
\(\order(\eta^4)\) in probability, before decoherence and model error are
added.  Table~\ref{tab:scaling} gives these powers with unit prefactors.  It
is a regime guide, not a predicted infidelity.

\begin{table}[ht]
\centering
\caption{Unit-prefactor scaling for \(\abs\alpha/2\pi=200\,\mathrm{MHz}\).
The entries are \(\eta^2\) and \(\eta^4\), not calculated gate errors.}
\label{tab:scaling}
\begin{tabular}{@{}rrrrr@{}}
\toprule
\(T\) (ns) & \(\abs\nu=\abs\alpha T\) & \(\eta\) & \(\order(\eta^2)\) &
\(\order(\eta^4)\) \\
\midrule
6.25 & 7.85 & 0.127 & \(1.62\times10^{-2}\) & \(2.63\times10^{-4}\) \\
10 & 12.57 & 0.0796 & \(6.33\times10^{-3}\) & \(4.01\times10^{-5}\) \\
20 & 25.13 & 0.0398 & \(1.58\times10^{-3}\) & \(2.51\times10^{-6}\) \\
40 & 50.27 & 0.0199 & \(3.96\times10^{-4}\) & \(1.57\times10^{-7}\) \\
\bottomrule
\end{tabular}
\end{table}

Published results show that the attainable improvement depends strongly on
duration and pulse family.  Motzoi~et~al.\ obtain reductions from roughly one
to several orders of magnitude in simulation, depending on the
waveform~\cite{Motzoi2009}.  Early transmon experiments demonstrated
derivative-of-Gaussian DRAG at few-nanosecond
duration~\cite{Chow2010}.  At \(4.16\,\mathrm{ns}\), closed-loop optimization
around DRAG reduced leakage from \(0.29\%\) to \(0.044\%\) and raised the
average Clifford fidelity from \(99.11\%\) to \(99.76\%\) on that
device~\cite{Werninghaus2021}.  Analytical FAST and higher-derivative DRAG
have subsequently reached leakage below \(3\times10^{-5}\) at
\(6.25\,\mathrm{ns}\), about a twentyfold reduction relative to a compared
cosine-DRAG pulse~\cite{Hyyppa2024}.  None of these values can be transferred
directly to a BARQ waveform: its spectrum, phase excursions, anharmonicity, and
coherence environment are different.  Experimental SCQC benchmarking also
identifies leakage as a gate-time-dependent limitation and selects duration by
balancing it against decoherence~\cite{SCQCExperiment2025}.

A numerical claim for this package needs the comparisons in
Table~\ref{tab:campaign}.  The minimum publication-quality comparison is a
qutrit simulation of the original BARQ pulse, the post-processed pulse, and a
pulse reoptimized with \eqref{eq:J}, all evaluated at the same peak amplitude
and coherence times.

\begin{table}[ht]
\centering
\caption{Simulation campaign needed before quoting an improvement.}
\label{tab:campaign}
{\small
\begin{tabular}{@{}>{\raggedright\arraybackslash}p{0.24\linewidth}>{\raggedright\arraybackslash}p{0.30\linewidth}>{\raggedright\arraybackslash}p{0.38\linewidth}@{}}
\toprule
Sweep & Quantities & Purpose \\
\midrule
\(T\) and \(\alpha\) & \(F_{\mathrm{avg}}\), \(L\), logical error &
  Locate leakage minima and the perturbative regime \\
\(\beta\) and Stark coefficient & two-dimensional error surface &
  Separate leakage cancellation from phase correction \\
Static \(\delta_z\) and gain & infidelity heat maps &
  Test whether robustness survives post-processing \\
Three to five levels & convergence of \(F_{\mathrm{avg}}\) and \(L\) &
  Detect higher-level and multiphoton leakage \\
\(T_1,T_\phi\) & open-system process fidelity &
  Choose a useful duration \\
Sampling and transfer function & filtered-pulse fidelity &
  Test sensitivity of derivative controls \\
\(\alpha\) and gain mismatch & worst-case and average fidelity &
  Quantify calibration robustness \\
\bottomrule
\end{tabular}}
\end{table}

\section{Conclusion}

DRAG fits SCQC as a correction in the pulse domain.  For a complex BARQ
envelope the first-order post-processor is \eqref{eq:drag-u}, supplemented by
a quadratic Stark detuning and by the endpoint conditions \(u=u'=0\).  This
exactly notches the bare envelope spectrum, but it only approximately preserves
the computational gate and the original closure condition.  The missing
robustness requirement is \eqref{eq:robust-blocks}, evaluated with the
corrected qutrit evolution; its first correction beyond ordinary curve closure
is the mixed integral \eqref{eq:mixed}.

Leakage-aware BARQ imposes the two complex constraints
\(\ell_0(\alpha T)=\ell_1(\alpha T)=0\) and can include their derivatives with
respect to uncertain parameters.  Exact qutrit propagation remains the
acceptance test, with \(L\) and \(F_{\mathrm{avg}}\) given by \eqref{eq:L}
and~\eqref{eq:F}.  Gate-time selection, endpoint smoothness, higher-derivative
filters, and spectral notches are compatible extensions.  A full
eight-dimensional \(\mathrm{SU}(3)\) construction is consistent with the
general geometric formalism, whereas the practical and geometrically faithful
intermediate step is closure of the perturbative leakage curve in
\(\mathbb{C}^2\).

\begin{thebibliography}{9}

\bibitem{Motzoi2009}
F.~Motzoi, J.~M.~Gambetta, P.~Rebentrost, and F.~K.~Wilhelm,
``Simple pulses for elimination of leakage in weakly nonlinear qubits,''
Phys.\ Rev.\ Lett.\ \textbf{103}, 110501 (2009),
\url{https://arxiv.org/abs/0901.0534}.

\bibitem{Chow2010}
J.~M.~Chow, L.~DiCarlo, J.~M.~Gambetta, F.~Motzoi, L.~Frunzio, S.~M.~Girvin,
and R.~J.~Schoelkopf,
``Optimized driving of superconducting artificial atoms for improved
single-qubit gates,''
Phys.\ Rev.\ A \textbf{82}, 040305 (2010).

\bibitem{Schutjens2013}
R.~Schutjens, F.~Abu~Dagga, D.~J.~Egger, and F.~K.~Wilhelm,
``Single-qubit gates in frequency-crowded transmon systems,''
Phys.\ Rev.\ A \textbf{88}, 052330 (2013),
\url{https://arxiv.org/abs/1306.2279}.

\bibitem{Vesterinen2014}
V.~Vesterinen, O.-P.~Saira, A.~Bruno, and L.~DiCarlo,
``Mitigating information leakage in a crowded spectrum of weakly anharmonic
qubits,''
arXiv:1405.0450 (2014).

\bibitem{Werninghaus2021}
M.~Werninghaus, D.~J.~Egger, F.~Roy, S.~Machnes, F.~K.~Wilhelm, and S.~Filipp,
``Leakage reduction in fast superconducting qubit gates via optimal control,''
npj Quantum Inf.\ \textbf{7}, 14 (2021),
\url{https://www.nature.com/articles/s41534-020-00346-2}.

\bibitem{Buterakos2021}
D.~Buterakos, S.~Das~Sarma, and E.~Barnes,
``Geometrical formalism for dynamically corrected gates in multiqubit
systems,''
PRX Quantum \textbf{2}, 010341 (2021).

\bibitem{BarnesReview}
E.~Barnes, F.~A.~Calderon-Vargas, W.~Dong, B.~Li, J.~Zeng, and F.~Zhuang,
``Dynamically corrected gates from geometric space curves,''
Quantum Sci.\ Technol.\ \textbf{7}, 023001 (2022),
\url{https://iopscience.iop.org/article/10.1088/2058-9565/ac4421}.

\bibitem{Hyyppa2024}
E.~Hyypp\"a \emph{et~al.},
``Reducing leakage of single-qubit gates for superconducting quantum processors
using analytical control pulse envelopes,''
PRX Quantum \textbf{5}, 030353 (2024).

\bibitem{SCQCExperiment2025}
``Implementing and benchmarking dynamically corrected gates on superconducting
devices using space curve quantum control,''
arXiv:2504.09767 (2025).

\end{thebibliography}

\end{document}
